\documentclass[a4paper,10pt]{ltjsarticle}
\usepackage[margin=23mm]{geometry}
\usepackage{booktabs}
\usepackage{longtable}
\usepackage{array}
\usepackage{amsmath}
\usepackage[hidelinks]{hyperref}

\title{dtq における多倍長融合積和 (\texttt{dw\_fma}/\texttt{tw\_fma}/\texttt{qw\_fma}) の\\
初等関数への適用 --- 性能と精度の評価}
\author{Tomonori Kouya}
\date{2026年8月10日}

\begin{document}
\maketitle

\begin{abstract}
多倍長浮動小数点演算ライブラリ dtq 0.0.3 で導入した分岐なし多倍長融合積和
(\texttt{dw\_fma}/\texttt{tw\_fma}/\texttt{qw\_fma}) が初等関数群の性能に与える効果を測定し,
さらに Taylor 級数・Horner 法の乗算--加算対を融合積和で置き換えた実験実装を評価した.
平方根を内部で用いる三角関数系は 0.0.2 比で最大 1.21 倍,Taylor/Horner の融合化により
\texttt{exp} 系は現行版比で最大 1.44 倍高速になった.MPFR (512\,bit) を基準とする
約 130 万サンプルの精度検証により,融合化による精度低下は実用上ないことを確認した.
併せて,既定ビルドでは \texttt{QD\_SLOPPY\_DIV} が優先されるため
\texttt{fma\_div} が使用されないという構成上の問題を報告する.
\end{abstract}

\section{はじめに}

dtq は double--double (\texttt{dd}), triple--double (\texttt{td}),
quad--double (\texttt{qd}) およびそれらの単精度版 (\texttt{ds}/\texttt{ts}/\texttt{qs})
を提供する多倍長浮動小数点演算ライブラリである.
バージョン 0.0.3 において,全 6 精度クラスに対する分岐なしの融合積和
$z = ab + c$ (\texttt{dw\_fma}/\texttt{tw\_fma}/\texttt{qw\_fma}) を導入した.
現行実装は ACS2026 の定式化に基づき,FPANVerifier と z3 により
記号精度 $p$ の下で誤差上界
\[ |z-(ab+c)| \le \{34u^2,\ 184u^3,\ 812u^4\}\,(|ab|+|c|) \quad (\text{それぞれ 2/3/4 語}) \]
と,全 FastTwoSum の前提条件,および出力の非重複性
(strongly dominates) が機械証明されている.

本レポートでは次の 3 点を評価する.
(1) 融合積和の導入 (0.0.2 $\to$ 0.0.3) が初等関数群の実行時間に与えた効果,
(2) 初等関数内部の Taylor 級数と \texttt{polyeval} (Horner 法) の
乗算--加算対を融合積和に置き換える追加最適化の効果,
(3) それらの精度への影響の網羅的検証.

\section{実験環境と測定方法}

\begin{itemize}
  \item CPU: Arm Cortex-X925/A725 (NVIDIA GB10 系,最大 3.9\,GHz),
        Linux 6.17.0-1021-nvidia (aarch64)
  \item コンパイラ: g++ 13.3.0.最適化フラグは全ビルド共通で
        \texttt{-O2 -ffp-contract=off} とし,configure 既定の
        \texttt{qd\_config.h} と \texttt{config.h} を共有した.
  \item 性能測定: 全語が非零のランダムなオペランド 512 個の配列に対して関数を適用し,
        1 パス 30\,ms 以上になるよう反復回数を較正した上で 7 回実行した最小値を採用
        (ns/回).実行コアは \texttt{taskset} で固定した.
        乱数種は固定であり,全ビルドが同一入力を処理する.
        結果の総和 (検算値) がビルド間で一致することを毎回確認した.
  \item 比較対象: (a) 0.0.2 (git \texttt{7edea02}),
        (b) 現行 main (\texttt{1989200},ACS2026 版 FMA),
        (c) 現行 main + Taylor/Horner 融合化 (実験実装,第\ref{sec:horner}節).
\end{itemize}

なお \texttt{td}/\texttt{ts} の初等関数はそれぞれ \texttt{qd}/\texttt{td} に昇格して
計算する実装のため,\texttt{td} の測定値は \texttt{qd} とほぼ一致する.
また \texttt{ds}/\texttt{ts}/\texttt{qs} の初等関数は倍精度クラスへの委譲であるため,
以下では \texttt{dd}/\texttt{td}/\texttt{qd} を対象とする.

\section{FMA カーネル単体の性能}

表~\ref{tab:kern1},\ref{tab:kern2} に \texttt{tests/fma\_bench}
(オペランド 4096 個,5 回中最速,ns/演算) の結果を示す.
「旧」は 0.0.2 相当のコード経路 (乗算+加算,長除算,\texttt{sqrt\_legacy}),
「新」は融合積和とそれに基づく除算・平方根である.
独立形 $z_i = a_i b_i + c_i$ はスループット律速,連鎖形 $s \gets a_i b_i + s$
はレイテンシ律速 (内積の内側ループ) を表す.

\begin{table}[htbp]
\centering
\caption{融合積和: 乗算+加算との比較 (ns/演算)}
\label{tab:kern1}
\begin{tabular}{l rrr rrr}
\toprule
 & \multicolumn{3}{c}{独立形} & \multicolumn{3}{c}{連鎖形} \\
\cmidrule(lr){2-4} \cmidrule(lr){5-7}
型 & 旧 & 新 & 倍率 & 旧 & 新 & 倍率 \\
\midrule
\texttt{dd} & 1.01 & 0.83 & 1.22 & 3.80 & 3.85 & 0.99 \\
\texttt{td} & 4.66 & 3.91 & 1.19 & 11.44 & 11.20 & 1.02 \\
\texttt{qd} & 17.34 & 7.57 & 2.29 & 17.74 & 13.34 & 1.33 \\
\texttt{ds} & 0.51 & 0.41 & 1.23 & 3.81 & 3.85 & 0.99 \\
\texttt{ts} & 9.53 & 3.86 & 2.47 & 14.05 & 9.02 & 1.56 \\
\texttt{qs} & 33.67 & 8.17 & 4.12 & 32.57 & 18.41 & 1.77 \\
\bottomrule
\end{tabular}
\end{table}

\begin{table}[htbp]
\centering
\caption{融合積和に基づく除算と平方根 (ns/演算)}
\label{tab:kern2}
\begin{tabular}{l rrr rrr}
\toprule
 & \multicolumn{3}{c}{除算} & \multicolumn{3}{c}{平方根} \\
\cmidrule(lr){2-4} \cmidrule(lr){5-7}
型 & 旧 & 新 & 倍率 & 旧 & 新 & 倍率 \\
\midrule
\texttt{dd} & 6.67 & 3.35 & 1.99 & 2.57 & 1.91 & 1.35 \\
\texttt{td} & 55.73 & 20.66 & 2.70 & 148.82 & 45.87 & 3.24 \\
\texttt{qd} & 154.71 & 48.48 & 3.19 & 316.56 & 119.46 & 2.65 \\
\texttt{ds} & 0.65 & 0.65 & 1.00 & 2.52 & 1.93 & 1.31 \\
\texttt{ts} & 32.68 & 9.87 & 3.31 & 76.07 & 39.87 & 1.91 \\
\texttt{qs} & 160.83 & 42.25 & 3.81 & 182.67 & 118.66 & 1.54 \\
\bottomrule
\end{tabular}
\end{table}

\section{初等関数への波及 (0.0.2 $\to$ 0.0.3)}
\label{sec:elem23}

表~\ref{tab:elem23} に既定構成 (\texttt{QD\_SLOPPY\_DIV} 有効) での
初等関数の実行時間を示す.高速化の源泉はほぼすべて FMA ベースの新しい
\texttt{sqrt} である.\texttt{sin}/\texttt{cos} は Taylor 部で
$\cos = \sqrt{1-\sin^2}$ を,\texttt{atan2} は $r=\sqrt{x^2+y^2}$ を
1 回ずつ呼ぶため,\texttt{qd} で 1.16--1.21 倍となる.
一方 \texttt{exp}/\texttt{log} 系は熱経路に平方根を含まないため不変である.

\begin{table}[htbp]
\centering
\caption{初等関数の実行時間 (ns/回,既定構成): 0.0.2 vs 0.0.3}
\label{tab:elem23}
\small
\begin{tabular}{l rrr rrr rrr}
\toprule
 & \multicolumn{3}{c}{\texttt{dd}} & \multicolumn{3}{c}{\texttt{td}} & \multicolumn{3}{c}{\texttt{qd}} \\
\cmidrule(lr){2-4} \cmidrule(lr){5-7} \cmidrule(lr){8-10}
関数 & 0.0.2 & 0.0.3 & 倍率 & 0.0.2 & 0.0.3 & 倍率 & 0.0.2 & 0.0.3 & 倍率 \\
\midrule
除算 & 4.4 & 4.4 & 1.00 & 67.8 & 68.2 & 0.99 & 179.2 & 179.0 & 1.00 \\
乗算 & 1.4 & 1.4 & 1.00 & 5.0 & 5.0 & 1.00 & 22.5 & 22.5 & 1.00 \\
\texttt{sqrt} & 10.5 & 7.5 & 1.40 & 239.6 & 94.8 & 2.53 & 502.8 & 238.9 & 2.10 \\
\texttt{log} & 194.9 & 194.8 & 1.00 & 6458.5 & 6453.5 & 1.00 & 6465.4 & 6461.2 & 1.00 \\
\texttt{log10} & 203.9 & 203.9 & 1.00 & 6652.2 & 6642.4 & 1.00 & 6653.5 & 6649.7 & 1.00 \\
\texttt{pow} & 382.8 & 382.9 & 1.00 & 8553.4 & 8550.2 & 1.00 & 8576.3 & 8561.2 & 1.00 \\
\texttt{nroot}$(x,3)$ & 80.2 & 80.3 & 1.00 & 612.6 & 612.6 & 1.00 & 1254.0 & 1254.6 & 1.00 \\
\texttt{exp} & 167.2 & 167.7 & 1.00 & 2069.0 & 2068.8 & 1.00 & 2066.2 & 2064.2 & 1.00 \\
\texttt{sin} & 163.8 & 160.7 & 1.02 & 1661.2 & 1397.4 & 1.19 & 1652.9 & 1388.0 & 1.19 \\
\texttt{cos} & 163.8 & 160.9 & 1.02 & 1669.2 & 1400.4 & 1.19 & 1660.8 & 1392.4 & 1.19 \\
\texttt{tan} & 182.2 & 179.1 & 1.02 & 1935.4 & 1668.9 & 1.16 & 1922.2 & 1656.9 & 1.16 \\
\texttt{atan} & 283.0 & 276.3 & 1.02 & 6981.2 & 5913.9 & 1.18 & 6966.7 & 5903.0 & 1.18 \\
\texttt{asin} & 315.9 & 305.2 & 1.04 & 7565.2 & 6241.3 & 1.21 & 7558.4 & 6238.4 & 1.21 \\
\texttt{acos} & 317.3 & 307.0 & 1.03 & 7571.1 & 6246.9 & 1.21 & 7568.6 & 6242.4 & 1.21 \\
\texttt{sinh} & 189.1 & 188.8 & 1.00 & 2294.5 & 2293.2 & 1.00 & 2287.5 & 2288.0 & 1.00 \\
\texttt{cosh} & 186.4 & 186.3 & 1.00 & 2282.0 & 2282.0 & 1.00 & 2280.2 & 2278.9 & 1.00 \\
\texttt{tanh} & 203.9 & 203.7 & 1.00 & 2515.2 & 2509.2 & 1.00 & 2511.7 & 2503.3 & 1.00 \\
\bottomrule
\end{tabular}
\end{table}

\paragraph{構成上の問題: \texttt{fma\_div} が既定で無効}
\texttt{operator/} のプリプロセッサ分岐は \texttt{\#ifdef QD\_SLOPPY\_DIV}
を最優先に評価し,configure は既定で \texttt{QD\_SLOPPY\_DIV} を定義する.
このため既定ビルドの除算は現在も \texttt{sloppy\_div} であり,
\texttt{fma\_div} は使用されない (表~\ref{tab:elem23} で除算が 1.00 倍である理由).
\texttt{QD\_SLOPPY\_DIV} を無効にした構成では表~\ref{tab:divacc} の通り
\texttt{fma\_div} は \texttt{accurate\_div} より大幅に速く,かつ高精度であるから,
分岐順の見直し (\texttt{fma\_div} を既定とする) を推奨する.

\begin{table}[htbp]
\centering
\caption{除算 (ns/回,\texttt{QD\_SLOPPY\_DIV} 無効構成): \texttt{accurate\_div} (0.0.2) vs \texttt{fma\_div}}
\label{tab:divacc}
\begin{tabular}{l rrr}
\toprule
型 & \texttt{accurate\_div} & \texttt{fma\_div} & 倍率 \\
\midrule
\texttt{dd} & 20.8 & 9.9 & 2.10 \\
\texttt{td} & 104.5 & 73.1 & 1.43 \\
\texttt{qd} & 249.4 & 182.3 & 1.37 \\
\bottomrule
\end{tabular}
\end{table}

\section{Taylor 級数・Horner 法の融合積和化}
\label{sec:horner}

初等関数の内部には「乗算の直後に加算」という形が多数現れる.
これらを融合積和 1 回に置き換えると,多倍長演算で支配的な正規化
(renormalization) の回数を削減できる.実験実装では次の 4 パターンを
\texttt{dd}/\texttt{qd} の実装 (\texttt{td} は \texttt{qd} 昇格で自動的に恩恵を受ける)
に適用した.

\begin{enumerate}
  \item \textbf{Horner 法} (\texttt{polyeval}):
        $r \gets r x + c_i$ を \texttt{fma}$(r, x, c_i)$ に置換.
  \item \textbf{Taylor 級数の項加算} (\texttt{exp}, \texttt{sin\_taylor},
        \texttt{cos\_taylor}, \texttt{sincos\_taylor}):
        $t \gets p \cdot \mathit{inv\_fact}_i;\ s \gets s + t$ を
        $s \gets \texttt{fma}(p, \mathit{inv\_fact}_i, s)$ に置換.
        収束判定は $t$ の代わりに先頭語の積
        $|p_0 \cdot (\mathit{inv\_fact}_i)_0|$ (相対誤差 $\sim 2^{-52}$ の近似)
        で行う.しきい値との大小判定には十分な精度である.
  \item \textbf{\texttt{exp} の 2 乗連鎖}:
        引数縮小の復元 $s \gets 2s + s^2$ (\texttt{qd} で 16 段,\texttt{dd} で 9 段)
        を $s \gets \texttt{fma}(s, s, 2s)$ に置換.
        1 段あたり正規化 2 回 (\texttt{sqr} と加算) が 1 回になるため寄与が最大.
  \item \textbf{\texttt{log} の Newton 反復}:
        $x \gets x + a e^{-x} - 1$ を $x \gets \texttt{fma}(a, e^{-x}, x) - 1$ に置換.
\end{enumerate}

表~\ref{tab:horner} に現行 main と実験実装の比較を示す.
\texttt{exp} は \texttt{qd} で 1.44 倍,パターン 3 の寄与が支配的である.
\texttt{log}/\texttt{pow} は内部の \texttt{exp} 呼び出しを通じて,
\texttt{sinh}/\texttt{cosh}/\texttt{tanh} も同様に 1.3--1.4 倍となる.
三角関数系はコストの主体が \texttt{sqrt} と引数縮小であり Taylor 部が
数項で打ち切られるため 3--5\,\%にとどまる.
8 次多項式の Horner 評価は \texttt{dd} で 1.44 倍である
(連鎖 $r$ に対する乗算と加算の正規化 2 回が 1 回になるため,
表~\ref{tab:kern1} の連鎖形よりも効果が大きい).

\begin{table}[htbp]
\centering
\caption{Taylor/Horner 融合化の効果 (ns/回,既定構成): 現行 main vs 実験実装}
\label{tab:horner}
\small
\begin{tabular}{l rrr rrr rrr}
\toprule
 & \multicolumn{3}{c}{\texttt{dd}} & \multicolumn{3}{c}{\texttt{td}} & \multicolumn{3}{c}{\texttt{qd}} \\
\cmidrule(lr){2-4} \cmidrule(lr){5-7} \cmidrule(lr){8-10}
関数 & 現行 & 融合 & 倍率 & 現行 & 融合 & 倍率 & 現行 & 融合 & 倍率 \\
\midrule
\texttt{exp} & 167.7 & 130.4 & 1.29 & 2068.3 & 1436.7 & 1.44 & 2065.2 & 1430.7 & 1.44 \\
\texttt{log} & 194.8 & 157.8 & 1.23 & 6457.9 & 4524.8 & 1.43 & 6461.1 & 4526.8 & 1.43 \\
\texttt{log10} & 203.8 & 166.7 & 1.22 & 6663.0 & 4715.3 & 1.41 & 6647.6 & 4718.3 & 1.41 \\
\texttt{pow} & 382.7 & 306.7 & 1.25 & 8585.4 & 5981.8 & 1.44 & 8561.7 & 6002.6 & 1.43 \\
\texttt{sinh} & 188.9 & 152.9 & 1.24 & 2293.7 & 1676.9 & 1.37 & 2287.0 & 1673.7 & 1.37 \\
\texttt{cosh} & 186.3 & 150.5 & 1.24 & 2280.8 & 1646.7 & 1.39 & 2278.9 & 1643.2 & 1.39 \\
\texttt{tanh} & 203.7 & 167.4 & 1.22 & 2508.5 & 1896.9 & 1.32 & 2502.0 & 1888.5 & 1.32 \\
\texttt{sin} & 160.7 & 154.7 & 1.04 & 1404.4 & 1333.3 & 1.05 & 1387.9 & 1324.8 & 1.05 \\
\texttt{cos} & 160.8 & 154.5 & 1.04 & 1402.7 & 1337.3 & 1.05 & 1392.1 & 1327.7 & 1.05 \\
\texttt{tan} & 179.1 & 172.7 & 1.04 & 1681.4 & 1602.9 & 1.05 & 1657.8 & 1589.8 & 1.04 \\
\texttt{atan} & 276.1 & 270.6 & 1.02 & 5941.5 & 5713.9 & 1.04 & 5903.6 & 5703.8 & 1.04 \\
\texttt{asin} & 305.4 & 299.7 & 1.02 & 6260.7 & 6034.9 & 1.04 & 6237.5 & 6034.4 & 1.03 \\
\texttt{acos} & 306.9 & 300.7 & 1.02 & 6265.1 & 6044.6 & 1.04 & 6243.8 & 6042.0 & 1.03 \\
\texttt{polyeval}(8次) & 67.4 & 46.9 & 1.44 & 219.1 & 149.9 & 1.46 & 430.7 & 327.7 & 1.31 \\
\bottomrule
\end{tabular}
\end{table}

実験実装は \texttt{qd\_test},\texttt{pslq\_test},\texttt{fma\_test} の
全テストに合格し,ベンチマーク全 54 項目の検算値は現行版と一致した.

\section{精度検証}
\label{sec:acc}

\subsection{方法}

MPFR 4.2.2 (仮数 512\,bit,\texttt{qd} の約 2.4 倍の精度) を基準とし,
0.0.2 / 現行 main / 実験実装の 3 ビルドに固定シードで同一入力を与えた.
誤差は各型自身の機械イプシロン
($\varepsilon_{\texttt{dd}}=2^{-104}$,
 $\varepsilon_{\texttt{td}}=2^{-158}$,
 $\varepsilon_{\texttt{qd}}=2^{-209}$) を単位とする相対誤差
$|y - \hat y| / (|\hat y|\,\varepsilon)$ で評価する.
テストセットは通常域に加え,実装変更が影響し得る境界を狙って設計した:
\texttt{exp} の引数が $m\ln 2$ 直近 (縮小後引数が微小になり打ち切り判定の
全域を掃く),\texttt{sinh}/\texttt{tanh} の分岐点 $0.05$ 近傍と微小引数,
\texttt{asin}/\texttt{acos} の $\pm1$ 近傍,\texttt{log} の $1$ 近傍,
$k\pi/2$ 近傍と巨大引数の三角関数,および激しい桁落ちを起こす
$(x-1)^9$ 展開係数の Horner 評価である.
1 型あたり 31 セット,計 93 セット・約 130 万サンプルを評価した.

\subsection{結果}

付録の表~\ref{tab:accfull} に全 93 セットの結果を示す.要点は次の通りである.

\begin{itemize}
  \item \textbf{初等関数は 3 ビルドで統計が一致}した.狙い撃ちしたストレスセット
        (表~\ref{tab:accsum}) を含め,最大誤差・RMS の差は最下位ビットの
        ノイズレベルにとどまる.打ち切り判定を先頭語近似に変えた影響は
        $m\ln 2$ 近傍セットでも検出されなかった.
  \item \textbf{唯一の実質差は悪条件の \texttt{polyeval}}:
        $(x-1)^9$ を根の近傍で評価するセットで \texttt{qd} の最大誤差が
        $4.6\times10^{36}\varepsilon \to 9.6\times10^{36}\varepsilon$ と約 2.1 倍になった
        (\texttt{dd} は $+8\%$,\texttt{td} は同一).この領域は現行版でも
        全サンプルの 99.9\,\%が $16\varepsilon$ を超える条件数支配の領域であり,
        誤差は両版とも条件数 $\times\,\varepsilon$ でスケールし,無補償 Horner 法の
        理論上界内にある.\texttt{qw\_fma} の証明済み誤差定数 ($812u^4$) が
        分離した乗算+加算の定数より大きいことと整合する挙動である.
        良条件のランダム係数セットでは $16\varepsilon$ 超えは 1 万中 1--2 件で両版同等である.
  \item 改善した項目もある (\texttt{qd} \texttt{pow} 最大 $5.34\to4.46\varepsilon$,
        \texttt{td} \texttt{polyeval} ランダム係数 $656\to542\varepsilon$,
        \texttt{dd} \texttt{exp} $[-1,1]$ $1.28\to1.10\varepsilon$).
\end{itemize}

\begin{table}[htbp]
\centering
\caption{ストレスセットにおける最大誤差 ($\varepsilon_{\texttt{qd}}$ 単位,\texttt{qd})}
\label{tab:accsum}
\begin{tabular}{ll rrr}
\toprule
関数 & セット & 0.0.2 & 現行 & 融合版 \\
\midrule
\texttt{exp} & $[-1,1]$ & 0.0657 & 0.0657 & 0.0722 \\
\texttt{exp} & $m\ln 2$ 近傍 & 4.11 & 4.11 & 4.11 \\
\texttt{log} & $1$ 近傍 & $1.5\times10^{15}$ & $1.5\times10^{15}$ & $2.0\times10^{15}$ \\
\texttt{sin} & $[-10,10]$ & 447 & 447 & 447 \\
\texttt{asin} & $\pm1$ 近傍 & 0.104 & 0.104 & 0.104 \\
\texttt{sinh} & $0.05$ 近傍 & 11.3 & 11.3 & 11.3 \\
\texttt{tanh} & 微小引数 & 0.973 & 1.07 & 1.07 \\
\texttt{pow} & 一様乱数 & 5.34 & 5.34 & 4.46 \\
\texttt{polyeval} & $(x-1)^9$ & $4.6\times10^{36}$ & $4.6\times10^{36}$ & $9.6\times10^{36}$ \\
\texttt{polyeval} & ランダム係数 & 60.3 & 60.3 & 78.9 \\
\bottomrule
\end{tabular}
\end{table}

\subsection{融合化と無関係の既存の弱点}

以下は 3 ビルドで完全に一致しており (0.0.2 から存在),今回の変更とは無関係である.

\begin{itemize}
  \item \texttt{dd} の \texttt{nroot}$(x,3)$ が最大 $520\varepsilon$
        (1 万中 3672 件が $16\varepsilon$ 超え).\texttt{td}/\texttt{qd} は
        $0.7\varepsilon$ 以下であり,\texttt{dd} の Newton 実装固有の弱さである.
  \item \texttt{dd} の \texttt{exp} は広域引数で最大 $111\varepsilon$,
        $m\ln 2$ 近傍で $70.9\varepsilon$.
  \item 極端な引数の \texttt{log}/\texttt{log10} ($|\log x|$ が大きい領域) では
        $e^{-x}$ の下位語がサブノーマル化して消えるため,
        \texttt{td}/\texttt{qd} で大きな $\varepsilon$ 単位誤差が生じる (フォーマットの限界).
  \item $k\pi/2$ 近傍・巨大引数の三角関数,$\pm1$ 近傍の \texttt{asin}/\texttt{acos},
        $1$ 近傍の \texttt{log} は,関数値の桁落ちにより相対誤差が条件数どおり増大する
        (引数縮小とキャンセレーションの原理的限界).
\end{itemize}

\section{まとめ}

\begin{enumerate}
  \item 融合積和に基づく \texttt{sqrt} により,三角関数系は 0.0.2 比で
        最大 1.21 倍 (\texttt{td}/\texttt{qd}) 高速になった.
  \item Taylor 級数・Horner 法・\texttt{exp} の 2 乗連鎖・\texttt{log} の
        Newton 反復を融合積和で書き直すことで,さらに \texttt{exp} 1.44 倍,
        \texttt{log} 1.43 倍,\texttt{sinh}/\texttt{cosh} 1.37--1.39 倍
        (\texttt{td}/\texttt{qd}),Horner 評価 1.31--1.46 倍を得た.
  \item 約 130 万サンプルの MPFR 検証により,融合化による精度低下は
        実用上ないことを確認した.唯一の注意点は悪条件多項式の Horner 評価で
        最大誤差が同一オーダー内で約 2 倍になり得ることである.
  \item 既定ビルドでは \texttt{QD\_SLOPPY\_DIV} が \texttt{fma\_div} に優先する.
        分岐順の是正を推奨する.
  \item 今後の課題: \texttt{dd} \texttt{nroot} の精度改善,
        実験実装の本体への取り込み.
\end{enumerate}

\appendix

\section{精度検証の全結果}

表~\ref{tab:accfull} に全 93 セットの最大誤差と RMS を示す
(各型の $\varepsilon$ 単位).$n$ はセットあたりのサンプル数で,
区間セットは 20{,}000,その他は 10{,}000 である.
``$\dagger$'' は 1 サンプルが \texttt{dd} 精度における $\tan$ の極に一致し
NaN を返したため RMS が定義されないことを表す (3 ビルドで同一サンプル).

\begingroup
\small
\setlength{\tabcolsep}{4pt}
\begin{longtable}{lll rrr rr}
\caption{全精度検証結果 (最大誤差と RMS,各型 $\varepsilon$ 単位)}
\label{tab:accfull} \\
\toprule
 & & & \multicolumn{3}{c}{最大誤差} & \multicolumn{2}{c}{RMS} \\
\cmidrule(lr){4-6} \cmidrule(lr){7-8}
型 & 関数 & セット & 0.0.2 & 現行 & 融合版 & 現行 & 融合版 \\
\midrule
\endfirsthead
\toprule
型 & 関数 & セット & \multicolumn{1}{c}{0.0.2} & 現行 & 融合版 & 現行 & 融合版 \\
\midrule
\endhead
\bottomrule
\endfoot
\texttt{dd} & \texttt{div} & 一様乱数 & 1.19 & 1.19 & 1.19 & 0.144 & 0.144 \\
\texttt{dd} & \texttt{sqrt} & 対数一様 & 1.44 & 1.34 & 1.34 & 0.25 & 0.25 \\
\texttt{dd} & \texttt{nroot3} & 対数一様 & 520 & 520 & 520 & 68.5 & 68.5 \\
\texttt{dd} & \texttt{exp} & $[-1,1]$ & 1.28 & 1.28 & 1.1 & 0.189 & 0.19 \\
\texttt{dd} & \texttt{exp} & $[-280,280]$ & 111 & 111 & 111 & 15 & 15 \\
\texttt{dd} & \texttt{exp} & $m\ln 2$ 近傍 & 70.9 & 70.9 & 70.9 & 11.4 & 11.4 \\
\texttt{dd} & \texttt{log} & 対数一様 & 62.5 & 62.5 & 62.5 & 13.7 & 13.7 \\
\texttt{dd} & \texttt{log} & $1$ 近傍 & $8.3\times10^{15}$ & $8.3\times10^{15}$ & $8.3\times10^{15}$ & $2.0\times10^{14}$ & $2.0\times10^{14}$ \\
\texttt{dd} & \texttt{log10} & 対数一様 & 62.9 & 62.9 & 62.9 & 13.6 & 13.6 \\
\texttt{dd} & \texttt{pow} & 一様乱数 & 40.9 & 40.9 & 41.6 & 5.74 & 5.89 \\
\texttt{dd} & \texttt{sin} & $[-10,10]$ & $3.6\times10^{3}$ & $3.6\times10^{3}$ & $3.6\times10^{3}$ & 33 & 33 \\
\texttt{dd} & \texttt{cos} & $[-10,10]$ & $4.6\times10^{3}$ & $4.6\times10^{3}$ & $4.6\times10^{3}$ & 48.7 & 48.7 \\
\texttt{dd} & \texttt{tan} & $[-10,10]$ & $4.6\times10^{3}$ & $4.6\times10^{3}$ & $4.6\times10^{3}$ & 58.8 & 58.8 \\
\texttt{dd} & \texttt{sin} & $k\pi/2$ 近傍 & $3.1\times10^{35}$ & $3.1\times10^{35}$ & $3.1\times10^{35}$ & $2.6\times10^{34}$ & $2.6\times10^{34}$ \\
\texttt{dd} & \texttt{cos} & $k\pi/2$ 近傍 & $3.1\times10^{35}$ & $3.1\times10^{35}$ & $3.1\times10^{35}$ & $2.6\times10^{34}$ & $2.6\times10^{34}$ \\
\texttt{dd} & \texttt{tan} & $k\pi/2$ 近傍 & $3.1\times10^{35}$ & $3.1\times10^{35}$ & $3.1\times10^{35}$ & $\dagger$ & $\dagger$ \\
\texttt{dd} & \texttt{sin} & 巨大引数 & $3.0\times10^{15}$ & $3.0\times10^{15}$ & $3.0\times10^{15}$ & $4.1\times10^{13}$ & $4.1\times10^{13}$ \\
\texttt{dd} & \texttt{cos} & 巨大引数 & $7.0\times10^{15}$ & $7.0\times10^{15}$ & $7.0\times10^{15}$ & $9.8\times10^{13}$ & $9.8\times10^{13}$ \\
\texttt{dd} & \texttt{asin} & $[-1,1]$ & 3.39 & 3.39 & 3.39 & 0.402 & 0.402 \\
\texttt{dd} & \texttt{acos} & $[-1,1]$ & 514 & 514 & 514 & 4.03 & 4.02 \\
\texttt{dd} & \texttt{asin} & $\pm1$ 近傍 & $3.4\times10^{6}$ & $3.4\times10^{6}$ & $3.4\times10^{6}$ & $2.0\times10^{5}$ & $2.0\times10^{5}$ \\
\texttt{dd} & \texttt{acos} & $\pm1$ 近傍 & $4.2\times10^{14}$ & $4.2\times10^{14}$ & $4.2\times10^{14}$ & $1.9\times10^{13}$ & $1.9\times10^{13}$ \\
\texttt{dd} & \texttt{atan} & 対数一様 & 1.68 & 1.79 & 1.79 & 0.143 & 0.144 \\
\texttt{dd} & \texttt{sinh} & $[-5,5]$ & 4.23 & 4.23 & 4.23 & 0.479 & 0.475 \\
\texttt{dd} & \texttt{cosh} & $[-5,5]$ & 2.1 & 2.1 & 2.1 & 0.357 & 0.356 \\
\texttt{dd} & \texttt{tanh} & $[-5,5]$ & 4.47 & 4.47 & 4.47 & 0.349 & 0.344 \\
\texttt{dd} & \texttt{sinh} & $0.05$ 近傍 & 8.79 & 8.79 & 8.79 & 0.897 & 0.901 \\
\texttt{dd} & \texttt{tanh} & $0.05$ 近傍 & 9.16 & 9.16 & 9.16 & 0.927 & 0.931 \\
\texttt{dd} & \texttt{sinh} & 微小引数 & 0.641 & 0.641 & 0.641 & 0.0329 & 0.0329 \\
\texttt{dd} & \texttt{tanh} & 微小引数 & 1.42 & 1.42 & 1.42 & 0.123 & 0.123 \\
\texttt{dd} & \texttt{polyeval} & $(x-1)^9$ & $2.5\times10^{37}$ & $2.5\times10^{37}$ & $2.7\times10^{37}$ & $1.2\times10^{36}$ & $1.1\times10^{36}$ \\
\texttt{dd} & \texttt{polyeval} & ランダム係数 & 653 & 653 & 743 & 9.19 & 10.1 \\
\texttt{td} & \texttt{div} & 一様乱数 & 0.748 & 0.748 & 0.748 & 0.0439 & 0.0439 \\
\texttt{td} & \texttt{sqrt} & 対数一様 & 0.309 & 0.319 & 0.319 & 0.0482 & 0.0482 \\
\texttt{td} & \texttt{nroot3} & 対数一様 & 0.677 & 0.677 & 0.677 & 0.0554 & 0.0554 \\
\texttt{td} & \texttt{exp} & $[-1,1]$ & 0.0305 & 0.0305 & 0.0305 & $7.5\times10^{-3}$ & $7.5\times10^{-3}$ \\
\texttt{td} & \texttt{exp} & $[-280,280]$ & 0.0312 & 0.0312 & 0.0312 & $7.5\times10^{-3}$ & $7.5\times10^{-3}$ \\
\texttt{td} & \texttt{exp} & $m\ln 2$ 近傍 & 0.0312 & 0.0312 & 0.0312 & $3.8\times10^{-3}$ & $3.8\times10^{-3}$ \\
\texttt{td} & \texttt{log} & 対数一様 & $4.2\times10^{10}$ & $4.2\times10^{10}$ & $4.2\times10^{10}$ & $6.7\times10^{8}$ & $6.7\times10^{8}$ \\
\texttt{td} & \texttt{log} & $1$ 近傍 & 0.19 & 0.19 & 0.232 & $4.6\times10^{-3}$ & $4.7\times10^{-3}$ \\
\texttt{td} & \texttt{log10} & 対数一様 & $3.9\times10^{10}$ & $3.9\times10^{10}$ & $3.9\times10^{10}$ & $6.3\times10^{8}$ & $6.3\times10^{8}$ \\
\texttt{td} & \texttt{pow} & 一様乱数 & 4.85 & 4.85 & 4.85 & 0.385 & 0.385 \\
\texttt{td} & \texttt{sin} & $[-10,10]$ & 0.0304 & 0.0304 & 0.0304 & $6.9\times10^{-3}$ & $6.9\times10^{-3}$ \\
\texttt{td} & \texttt{cos} & $[-10,10]$ & 0.03 & 0.03 & 0.03 & $6.7\times10^{-3}$ & $6.7\times10^{-3}$ \\
\texttt{td} & \texttt{tan} & $[-10,10]$ & 0.031 & 0.031 & 0.031 & $7.7\times10^{-3}$ & $7.7\times10^{-3}$ \\
\texttt{td} & \texttt{sin} & $k\pi/2$ 近傍 & $4.0\times10^{35}$ & $4.0\times10^{35}$ & $4.0\times10^{35}$ & $2.0\times10^{34}$ & $2.0\times10^{34}$ \\
\texttt{td} & \texttt{cos} & $k\pi/2$ 近傍 & $4.0\times10^{35}$ & $4.0\times10^{35}$ & $4.0\times10^{35}$ & $2.1\times10^{34}$ & $2.1\times10^{34}$ \\
\texttt{td} & \texttt{tan} & $k\pi/2$ 近傍 & $4.0\times10^{35}$ & $4.0\times10^{35}$ & $4.0\times10^{35}$ & $2.9\times10^{34}$ & $2.9\times10^{34}$ \\
\texttt{td} & \texttt{sin} & 巨大引数 & 0.0373 & 0.0373 & 0.0373 & $6.9\times10^{-3}$ & $6.9\times10^{-3}$ \\
\texttt{td} & \texttt{cos} & 巨大引数 & 0.0305 & 0.0305 & 0.0305 & $6.8\times10^{-3}$ & $6.8\times10^{-3}$ \\
\texttt{td} & \texttt{asin} & $[-1,1]$ & 0.0309 & 0.0309 & 0.0309 & $7.6\times10^{-3}$ & $7.6\times10^{-3}$ \\
\texttt{td} & \texttt{acos} & $[-1,1]$ & 0.0306 & 0.0306 & 0.0306 & $7.7\times10^{-3}$ & $7.7\times10^{-3}$ \\
\texttt{td} & \texttt{asin} & $\pm1$ 近傍 & 0.0251 & 0.0251 & 0.0251 & $7.5\times10^{-3}$ & $7.5\times10^{-3}$ \\
\texttt{td} & \texttt{acos} & $\pm1$ 近傍 & 0.031 & 0.031 & 0.031 & $7.5\times10^{-3}$ & $7.5\times10^{-3}$ \\
\texttt{td} & \texttt{atan} & 対数一様 & 0.0298 & 0.0298 & 0.0298 & $7.1\times10^{-3}$ & $7.1\times10^{-3}$ \\
\texttt{td} & \texttt{sinh} & $[-5,5]$ & 0.0298 & 0.0298 & 0.0298 & $7.6\times10^{-3}$ & $7.6\times10^{-3}$ \\
\texttt{td} & \texttt{cosh} & $[-5,5]$ & 0.031 & 0.031 & 0.031 & $7.8\times10^{-3}$ & $7.8\times10^{-3}$ \\
\texttt{td} & \texttt{tanh} & $[-5,5]$ & 0.0284 & 0.0284 & 0.0284 & $5.9\times10^{-3}$ & $5.9\times10^{-3}$ \\
\texttt{td} & \texttt{sinh} & $0.05$ 近傍 & 0.0308 & 0.0308 & 0.0308 & $6.2\times10^{-3}$ & $6.2\times10^{-3}$ \\
\texttt{td} & \texttt{tanh} & $0.05$ 近傍 & 0.0265 & 0.0265 & 0.0265 & $6.2\times10^{-3}$ & $6.2\times10^{-3}$ \\
\texttt{td} & \texttt{sinh} & 微小引数 & 0.03 & 0.03 & 0.03 & $2.9\times10^{-3}$ & $2.9\times10^{-3}$ \\
\texttt{td} & \texttt{tanh} & 微小引数 & 0.0299 & 0.0299 & 0.0299 & $3.0\times10^{-3}$ & $3.0\times10^{-3}$ \\
\texttt{td} & \texttt{polyeval} & $(x-1)^9$ & $1.0\times10^{37}$ & $1.0\times10^{37}$ & $1.0\times10^{37}$ & $3.4\times10^{35}$ & $3.7\times10^{35}$ \\
\texttt{td} & \texttt{polyeval} & ランダム係数 & 656 & 656 & 542 & 6.78 & 5.77 \\
\texttt{qd} & \texttt{div} & 一様乱数 & 1.12 & 1.12 & 1.12 & 0.0394 & 0.0394 \\
\texttt{qd} & \texttt{sqrt} & 対数一様 & 0.291 & 0.291 & 0.291 & 0.0301 & 0.0301 \\
\texttt{qd} & \texttt{nroot3} & 対数一様 & 0.7 & 0.7 & 0.7 & 0.0479 & 0.0479 \\
\texttt{qd} & \texttt{exp} & $[-1,1]$ & 0.0657 & 0.0657 & 0.0722 & $7.7\times10^{-3}$ & 0.0102 \\
\texttt{qd} & \texttt{exp} & $[-280,280]$ & 4.12 & 4.12 & 4.14 & 0.644 & 0.645 \\
\texttt{qd} & \texttt{exp} & $m\ln 2$ 近傍 & 4.11 & 4.11 & 4.11 & 0.61 & 0.61 \\
\texttt{qd} & \texttt{log} & 対数一様 & $4.4\times10^{26}$ & $4.4\times10^{26}$ & $4.4\times10^{26}$ & $5.9\times10^{24}$ & $5.9\times10^{24}$ \\
\texttt{qd} & \texttt{log} & $1$ 近傍 & $1.5\times10^{15}$ & $1.5\times10^{15}$ & $2.0\times10^{15}$ & $2.1\times10^{13}$ & $3.0\times10^{13}$ \\
\texttt{qd} & \texttt{log10} & 対数一様 & $3.4\times10^{26}$ & $3.4\times10^{26}$ & $3.4\times10^{26}$ & $4.2\times10^{24}$ & $4.2\times10^{24}$ \\
\texttt{qd} & \texttt{pow} & 一様乱数 & 5.34 & 5.34 & 4.46 & 0.413 & 0.522 \\
\texttt{qd} & \texttt{sin} & $[-10,10]$ & 447 & 447 & 447 & 3.87 & 3.87 \\
\texttt{qd} & \texttt{cos} & $[-10,10]$ & 43.2 & 43.2 & 43.2 & 0.544 & 0.544 \\
\texttt{qd} & \texttt{tan} & $[-10,10]$ & 447 & 447 & 447 & 3.91 & 3.91 \\
\texttt{qd} & \texttt{sin} & $k\pi/2$ 近傍 & $8.4\times10^{59}$ & $8.4\times10^{59}$ & $8.4\times10^{59}$ & $1.1\times10^{58}$ & $1.1\times10^{58}$ \\
\texttt{qd} & \texttt{cos} & $k\pi/2$ 近傍 & $9.2\times10^{59}$ & $9.2\times10^{59}$ & $9.2\times10^{59}$ & $1.4\times10^{58}$ & $1.4\times10^{58}$ \\
\texttt{qd} & \texttt{tan} & $k\pi/2$ 近傍 & $9.2\times10^{59}$ & $9.2\times10^{59}$ & $9.2\times10^{59}$ & $1.7\times10^{58}$ & $1.7\times10^{58}$ \\
\texttt{qd} & \texttt{sin} & 巨大引数 & $1.8\times10^{14}$ & $1.8\times10^{14}$ & $1.8\times10^{14}$ & $3.0\times10^{12}$ & $3.0\times10^{12}$ \\
\texttt{qd} & \texttt{cos} & 巨大引数 & $8.0\times10^{13}$ & $8.0\times10^{13}$ & $8.0\times10^{13}$ & $1.3\times10^{12}$ & $1.3\times10^{12}$ \\
\texttt{qd} & \texttt{asin} & $[-1,1]$ & 0.421 & 0.372 & 0.372 & 0.0489 & 0.0489 \\
\texttt{qd} & \texttt{acos} & $[-1,1]$ & 0.395 & 0.411 & 0.411 & 0.0279 & 0.0279 \\
\texttt{qd} & \texttt{asin} & $\pm1$ 近傍 & 0.104 & 0.104 & 0.104 & $3.6\times10^{-3}$ & $3.6\times10^{-3}$ \\
\texttt{qd} & \texttt{acos} & $\pm1$ 近傍 & 0.426 & 0.368 & 0.378 & 0.0199 & 0.0206 \\
\texttt{qd} & \texttt{atan} & 対数一様 & 1.39 & 1.39 & 1.36 & 0.0408 & 0.0413 \\
\texttt{qd} & \texttt{sinh} & $[-5,5]$ & 1.57 & 1.57 & 1.55 & 0.0565 & 0.0573 \\
\texttt{qd} & \texttt{cosh} & $[-5,5]$ & 0.66 & 0.66 & 0.66 & 0.0341 & 0.0346 \\
\texttt{qd} & \texttt{tanh} & $[-5,5]$ & 1.64 & 1.64 & 1.6 & 0.0725 & 0.0726 \\
\texttt{qd} & \texttt{sinh} & $0.05$ 近傍 & 11.3 & 11.3 & 11.3 & 0.272 & 0.273 \\
\texttt{qd} & \texttt{tanh} & $0.05$ 近傍 & 11.9 & 11.9 & 11.9 & 0.288 & 0.289 \\
\texttt{qd} & \texttt{sinh} & 微小引数 & 0.0934 & 0.0934 & 0.0934 & $3.6\times10^{-3}$ & $3.6\times10^{-3}$ \\
\texttt{qd} & \texttt{tanh} & 微小引数 & 0.973 & 1.07 & 1.07 & 0.04 & 0.04 \\
\texttt{qd} & \texttt{polyeval} & $(x-1)^9$ & $4.6\times10^{36}$ & $4.6\times10^{36}$ & $9.6\times10^{36}$ & $1.5\times10^{35}$ & $2.0\times10^{35}$ \\
\texttt{qd} & \texttt{polyeval} & ランダム係数 & 60.3 & 60.3 & 78.9 & 0.658 & 0.802 \\
\end{longtable}
\endgroup

\end{document}
